%%%% transaction-agent-ijcai26-preprint.tex
%%%% IJCAI--ECAI 26 preprint for the Transaction-Agent research prototype.

\typeout{IJCAI--ECAI 26 Instructions for Authors}

\documentclass{article}
\pdfpagewidth=8.5in
\pdfpageheight=11in

% The file ijcai26.sty is a copy from ijcai22.sty
% The file ijcai22.sty is NOT the same as previous years'
\usepackage{ijcai26}

% Use the postscript times font!
\usepackage{times}
\usepackage{soul}
\usepackage{url}
\usepackage[hidelinks]{hyperref}
\usepackage[utf8]{inputenc}
\usepackage[small]{caption}
\usepackage{graphicx}
\usepackage{amsmath}
\usepackage{amssymb}
\usepackage{amsthm}
\usepackage{booktabs}
\usepackage{algorithm}
\usepackage{algorithmic}
% Line numbers: the lineno package interferes with two-column float placement
% (tables/algorithm get overprinted by body text), so it is not loaded here.
% Re-enable with "\usepackage[switch]{lineno}" + "\linenumbers" ONLY if a
% blind submission system explicitly requires it, then re-verify the layout.
%\usepackage[switch]{lineno}

% Line numbers: enable only for anonymous submission. Keep disabled for the
% public preprint.
%\linenumbers

\urlstyle{same}

% the following package is optional:
%\usepackage{latexsym}

% See https://www.overleaf.com/learn/latex/theorems_and_proofs
% for a nice explanation of how to define new theorems, but keep
% in mind that the amsthm package is already included in this
% template and that you must *not* alter the styling.
\newtheorem{example}{Example}
\newtheorem{theorem}{Theorem}

% PDF Info Is REQUIRED.

% Please leave this \pdfinfo block untouched both for the submission and
% Camera Ready Copy. Do not include Title and Author information in the pdfinfo section
\pdfinfo{
/TemplateVersion (IJCAI.2026.0)
}

\title{Belief-Based Sequential Decision Making for Financial Transaction Fraud Detection}

\author{
    Rahul Krishna K$^1$\\
    \affiliations
    $^1$Independent Researcher\\
    \emails
    rahulkrishnak448@gmail.com
}

\begin{document}

\maketitle

\begin{abstract}
Financial transaction fraud detection is a sequential decision problem under
uncertainty: the true state of a transaction is hidden, the cost of a wrong
decision is asymmetric, and a cautious agent may defer an irreversible action
by pausing the transaction and collecting additional evidence. We present a
research prototype of a transaction-decision agent that (i)~summarizes a
customer's historical behavior into a lightweight profile, (ii)~extracts
rule-based evidence signals by comparing an incoming transaction to that
profile, (iii)~maintains an explicit Bayesian belief over the hidden
legitimate/fraudulent state, and (iv)~selects among \textsc{Approve},
\textsc{Hold}, and \textsc{Stop} by minimizing expected cost, where
\textsc{Hold} triggers a customer-verification feedback loop that updates the
belief before a final decision. We evaluate the agent on a frozen,
ground-truth-blind benchmark of 35 adversarially constructed cases. The agent
matches the expected decision on 22/35 cases (62.9\%). We decompose all 13
mismatches into three categories: evidence dilution under conditional
independence, decision-boundary calibration, and signals not yet modeled
(transaction velocity, merchant-category familiarity, profile maturity). Using
twin-case pairs, we empirically demonstrate that the agent's decisions do not
leak ground-truth labels. The system is implemented in pure Python (standard
library only; 23 passing unit and integration tests).
\end{abstract}

\section{Introduction}

Fraud detection for financial transactions is dominated by machine-learning
classifiers that map a transaction to a fraud score and then threshold the
score~\cite{bolton2002,bhattacharyya2011,leborgen2021}. Two properties of the
problem are frequently underspecified in this paradigm. First, the cost of a
mis-decision is strongly asymmetric and context-dependent: approving a
fraudulent high-value transfer is far more costly than delaying a legitimate
low-value purchase, and the two errors have different business
consequences~\cite{bahnsen2013}. Second, fraud detection is inherently a
\emph{sequential} decision problem: the transaction arrives with incomplete
information, and the agent may avoid an irreversible decision by pausing the
transaction and soliciting verification evidence from the customer.

This paper describes a research prototype --- the \emph{Transaction Agent} ---
that treats the problem as one of decision-making under uncertainty rather
than of classification. The agent (i)~builds a per-customer behavioral profile
from historical transactions, (ii)~extracts interpretable evidence signals by
comparing the incoming transaction with the profile, (iii)~updates an explicit
Bayesian belief over the hidden state using a small, fully transparent
likelihood table, and (iv)~chooses among \textsc{Approve}, \textsc{Hold}, and
\textsc{Stop} by minimizing expected cost. \textsc{Hold} is a first-class
information-gathering action: it pauses the transaction, requests customer
verification, and feeds the response back into the belief engine before a
final decision.

The prototype avoids all learned components: there are no classifiers, no
gradient updates, and no pretrained models. Every decision can be traced to a
handful of explicitly stated probabilities and costs, which makes the system
amenable to audit and to scientific study of its failure modes. We evaluate it
on a frozen benchmark of 35 adversarially constructed scenarios in which the
agent is kept strictly blind to ground truth. Our contributions are:

\begin{itemize}
  \item A complete, standard-library-only implementation of a sequential
  Bayesian decision agent for transaction fraud, with \textsc{Hold} as a
  first-class evidence-gathering action, validated by 23 automated tests.
  \item A ground-truth-blind evaluation protocol on a frozen 35-case benchmark,
  including \emph{twin-case} pairs with identical evidence but opposite hidden
  states, used to empirically verify the absence of label leakage.
  \item A quantitative failure taxonomy of the 13 benchmark mismatches,
  isolating three concrete modeling gaps: conditional-independence dilution,
  decision-boundary calibration, and missing signal channels (velocity,
  merchant-category familiarity, profile maturity).
\end{itemize}

\section{Related Work}

\emph{Statistical fraud detection.} Bolton and Hand~\cite{bolton2002} survey
statistical methods for fraud detection and emphasize that fraud detection is
a decision problem in which the cost structure of errors matters more than raw
accuracy. Hand and Henley~\cite{hand1997} reach a similar conclusion for
consumer credit scoring, treating scorecard thresholds as decision points that
should reflect the economics of lending. Our work follows this decision-theoretic
tradition rather than the accuracy-centric view of modern classifier benchmarks.

\emph{Cost-sensitive learning.} Bahnsen et al.~\cite{bahnsen2013} provide the
closest methodological ancestor: they formulate credit-card fraud detection as
a Bayes minimum-risk problem in which classification is performed by
minimizing the expected misclassification cost instead of the error rate. We
adopt the same principle but replace their expensive per-model cost matrix
estimation with an explicit, interpretable cost model and extend the single
decision into a \emph{sequential} loop in which one action
(\textsc{Hold}) yields additional evidence.

\emph{Modern fraud-detection benchmarks.} Large public data efforts such as the
IEEE-CIS Kaggle competition~\cite{iececis2019} and the reproducible
Fraud-Detection Handbook~\cite{leborgen2021} provide transactional data and
accurate/recalibrated classifier baselines. These datasets are valuable for
future validation of the agent's belief engine, but their evaluation protocol
is classification-oriented; they do not model deferred decisions or
information acquisition, which are central to our setting.

\emph{Uncertainty in agents.} The emerging literature on uncertainty
quantification in interactive LLM agents studies how an agent's estimate of
its own uncertainty changes as it interacts and gathers information. Our
verification loop is a concrete, non-learned instance of this principle: the
agent tracks a calibrated probability and chooses to gather more evidence only
when its belief falls in an ambiguous region. This connects the classical
decision-theoretic machinery with the emerging ``act-to-resolve-uncertainty''
view of agents.

\section{Problem Statement}

A transaction $T$ arrives with observable features: amount $A$, merchant $M$,
timestamp $\tau$, and location $L$. The \emph{hidden state} $H \in
\{\textsc{Legit}, \textsc{Fraud}\}$ is unobserved. The agent also has access to
the customer's historical transactions, from which a profile is derived. The
agent must select one of three actions:

\begin{itemize}
  \item \textsc{Approve}: release the transaction immediately;
  \item \textsc{Hold}: pause the transaction and request customer verification;
  \item \textsc{Stop}: terminate the transaction immediately.
\end{itemize}

Each action has a known cost $C(a \mid H)$ for each hidden state. The agent's
goal is to minimize the expected cost given its current belief. This is a
two-stage sequential decision: if \textsc{Hold} is chosen, a verification
response $V$ arrives and the belief is updated before the agent commits to a
terminal action. Because \textsc{Approve} and \textsc{Stop} are irreversible,
the problem is a degenerate (two-stage) partially observable Markov decision
process, in which the single ``act'' that changes the information state is the
verification solicitation.

\section{The Transaction Agent}

The agent's pipeline is summarized in Algorithm~\ref{alg:pipeline}. Its modules
are cleanly separated: profiling, evidence extraction, belief updating,
decision making, and verification. All parameters are explicit and modular;
each decision can be reproduced by hand from the tables in this paper.

\subsection{Behavioral Profiling}

Historical transactions for a customer $c$ are summarized into a profile
$S_c$:
\begin{itemize}
  \item \emph{Typical amount}: the median of all non-missing historical
  amounts;
  \item \emph{Common merchants}, \emph{typical locations}, and \emph{typical
  hours}: the sets of merchants, cities, and hour-of-day values appearing in
  the customer's history, in order of first appearance.
\end{itemize}

Missing values in the historical records are handled gracefully: a customer
with no usable amounts simply has no typical amount. In our experiments the
history consists of 200 synthetic records across 10 customers (C001--C010).

\subsection{Evidence Extraction}

Four rule-based checkers compare the incoming transaction with $S_c$ and emit
one evidence signal each:

\begin{itemize}
  \item \emph{Amount}: if $A > 1.5 \times \text{median}$, emit
  \texttt{large\_amount} (fraud-supporting); otherwise
  \texttt{amount\_consistent}.
  \item \emph{Merchant}: if $M \notin$ common merchants, emit
  \texttt{new\_merchant}; otherwise \texttt{merchant\_consistent}.
  \item \emph{Location}: if $L \notin$ typical locations, emit
  \texttt{unusual\_location}; otherwise \texttt{location\_consistent}.
  \item \emph{Time}: if $\tau.\text{hour} \notin$ typical hours, emit
  \texttt{unusual\_time}; otherwise \texttt{time\_consistent}.
\end{itemize}

Any missing feature emits a neutral \texttt{*\_unknown} signal that leaves the
belief unchanged. This keeps the system robust to incomplete transaction data.

\subsection{Bayesian Belief Updating}

The agent maintains a belief $(p_L, p_F)$ over the hidden state, initialized
with the prior
\begin{equation}
  p_L^{(0)} = 0.90, \qquad p_F^{(0)} = 0.10 .
\end{equation}
Each evidence signal $E$ updates the belief by Bayes' rule:
\begin{equation}
  p_F \leftarrow
  \frac{p_F \, P(E \mid \textsc{Fraud})}
  {p_F \, P(E \mid \textsc{Fraud}) + p_L \, P(E \mid \textsc{Legit})} .
\label{eq:bayes}
\end{equation}
Under the conditional-independence assumption, evidence arriving sequentially
is applied multiplicatively. In odds form the update is compact:
\begin{equation}
  O(F \mid E_1 \dots E_n) \;=\; O^{(0)}(F) \prod_{i=1}^{n} \Lambda(E_i),
\label{eq:odds}
\end{equation}
where each signal contributes a likelihood ratio
\begin{equation}
  \Lambda(E_i) = \frac{P(E_i \mid \textsc{Fraud})}
  {P(E_i \mid \textsc{Legit})},
\label{eq:lambda}
\end{equation}
and $O(F) = p_F / p_L$. Table~\ref{tab:likelihoods} lists the prototype
likelihood table, which is fully explicit and auditable. Consistency signals
carry likelihood ratios below 1 (they push the belief toward legitimate),
anomaly signals carry ratios above 1, and the two verification signals are the
strongest evidence available: a customer denial has ratio $19$, a confirmation
ratio $0.053$.

\subsection{Expected-Cost Decision Making}

The expected cost of action $a$ under belief $(p_L, p_F)$ is
\begin{equation}
  \mathbb{E}[C \mid a] = p_L \, C(a \mid \textsc{Legit})
  + p_F \, C(a \mid \textsc{Fraud}).
\label{eq:ec}
\end{equation}
The default cost model (Table~\ref{tab:costs}) sets approving a fraudulent
transaction at cost $10$, stopping a legitimate one at cost $8$, and treating
verification (\textsc{Hold}) as cheap and symmetric ($2$), reflecting low
customer friction for an out-of-band confirmation.

\begin{theorem}[Decision regions]
\label{thm:decision-regions}
Under the cost model of Table~\ref{tab:costs}, the expected-cost-minimizing
action is a function of $p_F$ alone:
\begin{equation}
  a^*(p_F) = \begin{cases}
    \textsc{Approve} & p_F < 0.2,\\
    \textsc{Hold}    & 0.2 \le p_F \le 0.75,\\
    \textsc{Stop}    & p_F > 0.75 .
  \end{cases}
\end{equation}
\end{theorem}

\begin{proof}
From Eq.~\eqref{eq:ec}: $\mathbb{E}[C \mid \textsc{Approve}] = 10 p_F$,
$\mathbb{E}[C \mid \textsc{Hold}] = 2$,
$\mathbb{E}[C \mid \textsc{Stop}] = 8(1 - p_F)$.
Comparing the first two, \textsc{Approve} is optimal iff $10 p_F < 2$,
i.e.\ $p_F < 0.2$. Comparing the last two, \textsc{Hold} beats \textsc{Stop}
iff $2 < 8(1 - p_F)$, i.e.\ $p_F < 0.75$.
\end{proof}

\begin{algorithm}[tb]
\caption{Transaction decision loop}
\label{alg:pipeline}
\textbf{Input}: Customer history $\mathcal{H}_c$, incoming transaction $T$\\
\textbf{Parameter}: Prior $(0.90, 0.10)$; LRs (Tab.~\ref{tab:likelihoods}); costs (Tab.~\ref{tab:costs})\\
\textbf{Output}: Final action $a^* \in \{\textsc{Approve}, \textsc{Stop}\}$
\begin{algorithmic}[1]
\STATE Build profile $S_c$ from $\mathcal{H}_c$ (median amount, merchants, hours).
\STATE Extract evidence $\{E_1,\dots,E_4\}$ by comparing $T$ with $S_c$.
\STATE Initialize belief $b \leftarrow (0.90, 0.10)$.
\FOR{each $E_i$}
  \STATE Update $b$ with Bayes' rule (Eq.~\eqref{eq:bayes}).
\ENDFOR
\STATE Compute $\mathbb{E}[C \mid a]$ for all $a$ (Eq.~\eqref{eq:ec}).
\IF{$\arg\min_a \mathbb{E}[C\mid a] = \textsc{Hold}$}
  \STATE Request customer verification; obtain response $V$.
  \STATE Convert $V$ to verification evidence (\texttt{customer\_confirms} /
  \texttt{customer\_denies}).
  \STATE Update $b$ again; recompute expected costs.
  \STATE Return the new argmin (over \textsc{Approve} / \textsc{Stop}).
\ELSE
  \STATE Return the argmin action directly.
\ENDIF
\end{algorithmic}
\end{algorithm}

\begin{table}[tb]
\centering
\footnotesize
\setlength{\tabcolsep}{4pt}
\begin{tabular}{lccc}
\toprule
Evidence signal & $P(E\mid\textsc{Frg})$ & $P(E\mid\textsc{Lgt})$ & $\Lambda$ \\
\midrule
\texttt{large\_amount}      & 0.30 & 0.20 & 1.50 \\
\texttt{amount\_consistent} & 0.70 & 0.80 & 0.875 \\
\texttt{new\_merchant}      & 0.30 & 0.15 & 2.00 \\
\texttt{merchant\_consistent}& 0.70 & 0.85 & 0.824 \\
\texttt{unusual\_location}  & 0.30 & 0.10 & 3.00 \\
\texttt{location\_consistent}& 0.70 & 0.90 & 0.778 \\
\texttt{unusual\_time}      & 0.40 & 0.10 & 4.00 \\
\texttt{time\_consistent}   & 0.60 & 0.90 & 0.667 \\
\texttt{customer\_confirms} & 0.05 & 0.95 & 0.053 \\
\texttt{customer\_denies}   & 0.95 & 0.05 & 19.00 \\
\bottomrule
\end{tabular}
\caption{Prototype likelihood table with ratios $\Lambda =
P(E\mid\textsc{Fraud})/P(E\mid\textsc{Legit})$. Consistency signals are
complements of anomaly signals; ``unknown'' signals are neutral and omitted.}
\label{tab:likelihoods}
\end{table}

\begin{table}[tb]
\centering
\footnotesize
\setlength{\tabcolsep}{4pt}
\begin{tabular}{lccc}
\toprule
Hidden state & \textsc{Approve} & \textsc{Hold} & \textsc{Stop} \\
\midrule
\textsc{Legit}  & 0.0 & 2.0 & 8.0 \\
\textsc{Fraud}  & 10.0 & 2.0 & 0.0 \\
\bottomrule
\end{tabular}
\caption{Default cost model. Asymmetric costs make the optimal action a pure
function of $p_F$ (Theorem~\ref{thm:decision-regions}).}
\label{tab:costs}
\end{table}

A subtle discovery from the prototype's development is worth reporting here:
the agent's \emph{first} cost model set $\textsc{Hold}$ at
$(\text{legit}=6,\ \text{fraud}=4)$, which makes \textsc{Hold} dominated --- its
expected cost exceeds both \textsc{Approve} and \textsc{Stop} at every belief
point in $[0,1]$. \textsc{Hold} was mathematically inadmissible, a fact that
was invisible in an accuracy-framed classifier but immediately surfaced in the
expected-cost formulation. The default model of Table~\ref{tab:costs} was
adopted specifically so that the three decision regions of
Theorem~\ref{thm:decision-regions} are non-degenerate.

\subsection{Verification as Information Gathering}

If the initial decision is \textsc{Hold}, the agent creates a pending
verification request, records the customer's response, converts it into
verification evidence (\texttt{customer\_confirms} or \texttt{customer\_denies}),
updates the belief once more, and re-runs expected-cost minimization over the
two terminal actions. Because a denial carries likelihood ratio 19, a single
failed verification can push a $p_F$ in the ambiguous $[0.2, 0.75]$ band above
the 75\% \textsc{Stop} threshold; a confirmation pushes it far below 20\%.
This realizes, in explicit probabilistic form, the principle that a
\emph{reversible} action at a 50/50 belief is often the cheapest choice.

\section{Evaluation}

\subsection{Benchmark Design and Ground-Truth Blindness}

The benchmark file used throughout this evaluation%
\footnote{\texttt{transaction\_agent\_test\_cases\_v2.csv} in \texttt{experiments/}.} is
\emph{frozen} and read-only:
it contains 35 scenario
cases (TC01--TC35) and must not be regenerated. Each case specifies the
observed features
(amount, merchant, time, location), a natural-language description of the
customer's historical behavior, explicit evidence notes, and --- hidden from
the agent --- the ground-truth state ($22$ legitimate, $13$ fraudulent) and the
human-expected action ($12$ \textsc{Approve}, $15$ \textsc{Hold},
$8$ \textsc{Stop}).

The evaluation runner computes the agent's decision from feature columns only.
Ground truth is read solely for post-hoc comparison, after the decision is
finalized. The set deliberately contains five \emph{twin pairs} --- cases with
identical evidence but opposite hidden states (e.g., TC04/TC05, TC10/TC11) ---
so that any leakage of ground truth into the decision would be directly
observable as a disagreement between twins.

\subsection{Results}

The agent selected \textsc{Approve} 21 times, \textsc{Hold} 9 times, and
\textsc{Stop} 5 times, matching the expected action on 22/35 cases
($62.9\%$). Table~\ref{tab:results} shows representative cases. Consistent
transactions (TC01, TC19) are approved with posterior fraud probability near
$4\%$; fully anomalous transactions (TC02, TC23, TC33) are stopped with
$p_F = 80\%$; and genuinely ambiguous two-anomaly cases (TC17, TC18) land in
the \textsc{Hold} band, exactly where the cost model intends.

\begin{table}[tb]
\centering
\footnotesize
\setlength{\tabcolsep}{2.5pt}
\begin{tabular}{llrcl}
\toprule
Case & Evidence summary & $p_F$ & Agt./Exp. & Match\\
\midrule
TC01 & 4 consistent & 0.040 & A/A & \checkmark\\
TC02 & 4 anomalies & 0.800 & S/S & \checkmark\\
TC04 & large amt + 3 ok & 0.066 & A/H & $\times$\\
TC07 & unusual time + 3 ok & 0.199 & A/H & $\times$\\
TC08 & 3 anom. + 1 ok & 0.509 & H/S & $\times$\\
TC17 & 2 anom. (time+loc) & 0.490 & H/H & \checkmark\\
TC20 & 3 txs in 6~min & 0.092 & A/S & $\times$\\
TC25 & unfamiliar merc. (cat.) & 0.092 & A/H & $\times$\\
TC33 & extreme amt + 4 anom. & 0.800 & S/S & \checkmark\\
TC35 & sparse profile ($<$5 txs) & 0.092 & A/H & $\times$\\
\bottomrule
\end{tabular}
\caption{Representative benchmark cases. $p_F$ is the posterior fraud
probability after belief updating; A/H/S denote
\textsc{Approve}/\textsc{Hold}/\textsc{Stop}; Match compares the agent's
decision with the human-expected action.}
\label{tab:results}
\end{table}

\subsection{Failure Taxonomy}

All 13 mismatches fall into three categories, each with a distinct mechanism:

\paragraph{Category A: Conditional-independence dilution (7 cases).}
TC04, TC05, TC10, TC11, TC13, TC22, TC07. A single moderate anomaly (e.g.\ a
large amount) combined with three routine consistent signals is diluted by the
multiplicative independence update. In TC04/TC05 the product of likelihood
ratios is $1.5 \times 0.824 \times 0.778 \times 0.667 = 0.642$, driving $p_F$
\emph{down} from 10\% to 6.6\% below the 20\% \textsc{Hold} threshold, even
though the amount is $4\times$ the historical maximum. TC07 reaches
$p_F = 19.94\%$, missing the boundary by 0.06 points. This is a modeling
limitation, not a software bug: naive conditional independence lets many
``routine'' signals cancel an extreme outlier.

\paragraph{Category B: Decision-boundary behavior (2 cases).}
TC08, TC14. The benchmark expected \textsc{Stop}; the agent chose
\textsc{Hold} with $p_F \in [0.40, 0.51]$. Because \textsc{Stop} costs $8$
against a legitimate transaction, its expected cost is $8(1-p_F) \ge 3.93$,
which exceeds the constant cost $2$ of \textsc{Hold}. Under the stated cost
model, pausing a 50/50 transaction for verification is strictly cheaper than
an irreversible cancellation --- the agent's choice is the rational one under
Theorem~\ref{thm:decision-regions}. The disagreement reveals a normative difference in
how aggressively fraud should be terminated, i.e.\ a cost-calibration question,
not a reasoning error.

\paragraph{Category C: Missing signal channels (4 cases).}
TC20 (three transactions within six minutes --- transaction \emph{velocity} is
not represented in the four single-transaction checkers), TC25/TC26 (an
unfamiliar merchant within a familiar spending category --- merchant familiarity
is currently binary), and TC35 (a profile with fewer than five prior
transactions --- profile \emph{maturity} is not modeled). Each of these is a
deliberately deferred feature whose signal channel does not yet exist.

\subsection{Twin-Pair Probe for Label Leakage}

For every twin pair (TC04/TC05, TC10/TC11, TC17/TC18, TC25/TC26, TC31/TC32),
the agent produced the identical decision for both members, despite opposite
ground-truth states. Since the two members share evidence and differ only in
the hidden label, identical outputs are consistent with the agent operating
purely on features. This serves as an empirical control that the ground-truth
columns do not influence the decision.

\section{Discussion}

The prototype's value is primarily investigative. Three conclusions emerge.
First, an explicit belief state makes otherwise subtle modeling assumptions
visible. The conditional-independence assumption is attractive for its
simplicity and its closed-form updates (Eq.~\eqref{eq:odds}), but the benchmark
shows it over-weights repeated routine evidence against a single extreme
anomaly. Second, expected-cost decision making surfaces normative tension that
accuracy metrics hide: the TC08/TC14 disagreements are not errors but a
stance on the economics of false stops, and the history of the cost model shows
how easily a classifier designer can make a reasonable action
(\textsc{Hold}) structurally unused. Third, the twin pairs and the frozen
benchmark provide a reusable evaluation protocol for future signal channels.

The system has deliberate limitations. All probabilities, likelihoods, and
costs are synthetic research assumptions, not calibrated banking statistics;
parameters must be estimated from real data before deployment. The evidence
set is small (four channels) and single-transaction; and the verification loop
assumes a reliable response channel. Deployment concerns (timing, fraud rings,
identity-verification cost) are out of scope for this prototype.

\section{Future Work}

The benchmark failures define the research agenda directly. (i)~\emph{Non-linear
evidence combination}: extreme individual anomalies should resist dilution,
e.g.\ by capping the strength of routine signals or introducing a generalized
aggregation rule; TC07's 0.06-point miss indicates that even the boundary
behavior is fragile. (ii)~\emph{New signal channels}: transaction velocity,
merchant-category familiarity, and profile-maturity estimates should become
first-class evidence channels with their own likelihood entries. (iii)~\emph{Cost
calibration}: the HOLD/STOP boundary should be studied against real
false-positive and false-negative economics, and the cost model should be
parameterized per transaction context (e.g., low versus high denomination).
(iv)~\emph{Empirical validation}: the agent's belief engine and likelihood
table should be confronted with real transaction data such as the IEEE-CIS
dataset~\cite{iececis2019} after estimating parameters from observed
frequencies. We also plan to study the connection between the verification
loop and uncertainty-driven information acquisition in agents.

\section*{Ethical Statement}

This paper reports a research prototype evaluated only on synthetic data.
No real customer data was used, and no deployment or financial decisions were
influenced by the system. The benchmark cases are fabricated and carry no
personally identifiable information.

\bibliographystyle{named}
\begin{thebibliography}{9}

\bibitem{bolton2002}
R.~J. Bolton and D.~J. Hand.
Statistical fraud detection: A review.
\emph{Statistical Science}, 17(3):235--255, 2002.

\bibitem{hand1997}
D.~J. Hand and W.~E. Henley.
Statistical classification methods in consumer credit scoring: A review.
\emph{Journal of the Royal Statistical Society: Series A}, 160(3):523--541, 1997.

\bibitem{bahnsen2013}
A.~C. Bahnsen, D.~Aouada, A.~Stojanovic, and B.~Ottersten.
Cost sensitive credit card fraud detection using Bayes minimum risk.
In \emph{Proc.~12th International Conference on Machine Learning and
Applications (ICMLA)}, volume~1, pages 333--338. IEEE, 2013.

\bibitem{bhattacharyya2011}
S.~Bhattacharyya, S.~Jha, K.~Tharakunnel, and J.~C. Westland.
Data mining for credit card fraud: A comparative study.
\emph{Decision Support Systems}, 50(3):602--613, 2011.

\bibitem{leborgen2021}
Y.-A.~Le~Borgne, W.~Siblini, B.~Lebichot, and G.~Bontempi.
Reproducible Machine Learning for Credit Card Fraud Detection ---
Practical Handbook.
Open-source handbook, 2021. \url{https://github.com/Fraud-Detection-Handbook}.

\bibitem{iececis2019}
IEEE-CIS.
IEEE-CIS Fraud Detection.
Kaggle competition data set, 2019.
\url{https://www.kaggle.com/c/ieee-fraud-detection}.

\end{thebibliography}

\end{document}